\SourceLeaf{166}%
% A544-P0166-D01
hingegen in dem Falle, wo $n$ von der Form $4m+3$ ist, die Summe der ersten Reihe $=0$, und die der zweiten $=\pm\surd n$ wird. Das der Wurzelgrösse vorzusetzende Zeichen hängt von dem Werthe der Zahl $k$ oder vielmehr von dessen Relation zu $n$ ab, und lässt sich leicht für \emph{alle} Werthe von $k$ bei einem gegebenen Werthe von $n$ bestimmen, sobald es für \emph{einen} bestimmt ist. Man kann nemlich zeigen, dass für alle Werthe von $k$, welche quadratische Reste von $n$ sind, durchaus \emph{einerlei} Zeichen gilt, und dann das entgegengesetzte für alle diejenigen, die quadratische Nichtreste von $n$ sind. Da in dem angeführten Werke die Untersuchung so weit bereits geführt, und nur die Bestimmung des Zeichens für irgend einen Werth von $k$ noch übrig war: so hätte man glauben sollen, dass nach Beseitigung der Hauptsache diese nähere Bestimmung sich leicht würde ergänzen lassen, um so mehr, da die Induction dafür sogleich ein äusserst einfaches Resultat gibt: für $k=1$, oder für alle Werthe, welche quadratische Reste von $n$ sind, muss nemlich die Wurzelgrösse in obigen Formeln durchaus \emph{positiv} genommen werden. Allein bei der Aufsuchung des Beweises dieser Bemerkung treffen wir auf ganz unerwartete Schwierigkeiten, und dasjenige Verfahren, welches so genugthuend zu der Bestimmung des absoluten Werths jener Reihen führte, wird durchaus unzureichend befunden, wenn es die vollständige Bestimmung der Zeichen gilt. Den \emph{metaphysischen} Grund dieses Phänomens (um den bei den Französischen Geometern üblichen Ausdruck zu gebrauchen) hat man in dem Umstande zu suchen, dass die Analyse bei der Theilung des Kreises zwischen den Bögen $\omega,2\omega,3\omega\ldots(n-1)\omega$ keinen Unterschied macht, sondern alle auf gleiche Art umfasst; und da hiedurch die Untersuchung ein neues Interesse erhält: so fand Hr. Prof. \textsc{Gauss} hierin gleichsam eine Aufforderung, nichts unversucht zu lassen, um die Schwierigkeit zu besiegen. Erst nach vielen und mannigfaltigen vergeblichen Versuchen ist ihm dieses auf einem auch an sich selbst merkwürdigen Wege gelungen. Er geht nemlich von der Summation einiger Reihen aus, deren Glieder unter folgender Form begriffen sind:
\[
\frac{(1-x^m)(1-x^{m-1})(1-x^{m-2})\ldots(1-x^{m-\mu+1})}{(1-x)\ (1-xx)\ (1-x^3)\dotsfour(1-x^\mu)}
\]
Bezeichnet man, der Kürze halber, eine solche Function durch $(m,\mu)$, welche, wie in der Abhandlung gezeigt wird, immer eine \emph{ganze} Function von $x$ ist: so brechen die Reihen
\SourceLeaf{167}%
% A544-P0167-D01
\[
1-(m,1)+(m,2)-(m,3)+\text{etc.}
\]
\[
1+x^{\frac12}(m,1)+x(m,2)+x^{\frac32}(m,3)+\text{etc.}
\]
nach dem $m+1^{\mathrm{sten}}$ Gliede ab, insofern $m$ eine ganze positive Zahl bedeutet, und die Summe der ersten Reihe wird für gerade Werthe von $m$
\[
=(1-x)(1-x^3)(1-x^5)\ldots(1-x^{m-1})
\]
und $=0$ für ungerade Werthe von $m$; hingegen die Summe der zweiten Reihe wird allemal
\[
=(1+x^{\frac12})(1+x)(1+x^{\frac32})\ldots(1+x^{\frac12m})
\]
Auch für gebrochene und negative Werthe von $m$ führt die Summation dieser Reihen auf interessante Resultate, obwohl dieselben zu der gegenwärtigen Absicht nicht nöthig sind: wir begnügen uns, nur eines derselben hier anzuführen. Die unendliche Reihe
\[
1+x+x^3+x^6+x^{10}+\text{etc.}
\]
wo die Exponenten die Trigonalzahlen sind, ist das Product aus den Factoren
\[
\frac{1-xx}{1-x}\times\frac{1-x^4}{1-x^3}\times\frac{1-x^6}{1-x^5}\times\frac{1-x^8}{1-x^7}\ \text{etc.}
\]
oder, wenn man lieber will, aus
\[
(1+x)^2(1+xx)^2(1+x^3)^2(1+x^4)^2\ \text{etc.}
\]
in
\[
(1-x)(1-xx)1-x^3)(1-x^4)\ \text{etc.}
\]

% A544-P0167-D02
Die Entwicklung der Art, wie diese Summationen auf den Hauptgegenstand angewandt werden, würde uns hier zu weit führen: wir dürfen die Leser um so eher auf diese selbst verweisen, da sie bald im Druck erscheinen wird. Jene oben angeführten Summationen sind nur eine specielle Anwendung von der Summation folgender Reihen:
\[
1+\cos k\omega+\cos4k\omega+\cos9k\omega+\text{etc.}+\cos(n-1)^2k\omega=T
\]
\[
\sin k\omega+\sin4k\omega+\sin9k\omega+\text{etc.}+\sin(n-1)^2k\omega=U
\]
welche in der Abhandlung für alle Werthe von $k$, und ohne die Einschränkung,
\SourceLeaf{168}%
% A544-P0168-D01
dass $n$ eine Primzahl sei, gelehrt wird. Es wird nemlich gezeigt, dass
\[
T=\pm\surd n,\quad T=\pm\surd n,\quad T=0,\quad T=0
\]
und
\[
U=\pm\surd n,\quad U=0,\quad U=0,\quad U=\pm\surd n
\]
wird, je nachdem $n$ von der Form $4m,4m+1,4m+2,4m+3$ resp. ist; das Zeichen der Wurzelgrösse hängt hier wiederum von $k$ ab, und die Unterscheidung vieler einzelner Fälle nöthig machende Bestimmung desselben auf zwei verschiedenen Wegen wird so entwickelt und bewiesen, dass nichts zu wünschen übrig bleiben wird. Die Vergleichung dieser beiden Wege unter sich führt noch auf folgenden sehr merkwürdigen Lehrsatz: Wenn $n$ das Product aus einer beliebigen Anzahl ungleicher ungerader Primzahlen $a,b,c,d$ u. s. w. ist, unter welchen sich zusammen $\mu$ von der Form $4m+3$ befinden: wenn ferner unter jenen Factoren zusammen $\nu$ vorkommen, von deren jedem das Product der übrigen (also resp. $\frac na,\frac nb,\frac nc,\frac nd$ u. s. w.) ein quadratischer Nichtrest ist; so wird $\nu$ gerade sein, so oft $\mu$ von der Form $4m$ oder $4m+1$ ist, hingegen ungerade, so oft $\mu$ von der Form $4m+2$ oder $4m+3$ ist. Von diesem Lehrsatze ist das bekannte Fundamental-Theorem bei den quadratischen Resten nur ein specieller Fall, sowie umgekehrt jener leicht aus diesem abgeleitet werden kann. Man sieht sich also durch diese Untersuchungen zugleich im Besitz von einem vierten Beweise dieses wichtigen Theorems, welches von dem Verf. zuerst auf zwei ganz verschiedenen Wegen in den \emph{Disquisitionibus Arithmeticis} und auf einem dritten eben so verschiedenen unlängst in einer eigenen Abhandlung bewiesen war.
\BookRule{.28\linewidth}

\clearpage
\SourceLeaf{169}%
% A544-P0169-D01
\pdfbookmark[0]{Göttingische gelehrte Anzeigen. 1817 März 10.}{notice1817march10}
\begin{center}
\rule{.82\linewidth}{.35pt}\par
Göttingische gelehrte Anzeigen.\quad 1817 März 10.\par
\rule{.82\linewidth}{.35pt}
\end{center}
Am 10. Februar wurde der Königl. Societät von Hrn. Hofr. \textsc{Gauss} eine Vorlesung eingereicht, überschrieben:
\begin{center}\itshape
Theorematis fundamentalis in doctrina de residuis quadraticis demonstrationes et ampliationes novae.
\end{center}

% A544-P0169-D02
Es ist eine Eigenthümlichkeit der höhern Arithmetik, dass so viele ihrer schönsten Lehrsätze mit grösster Leichtigkeit durch Induction entdeckt werden können, deren Beweise jedoch nichts weniger als nahe liegen, sondern oft erst nach vielen vergeblichen Versuchen mit Hülfe tiefeindringender Untersuchungen und glücklicher Combinationen gefunden werden. Dies merkwürdige Phänomen entspringt aus der oft wunderbaren Verkettung der verschiedenartigen Lehren in jenem Theile der Mathematik, und eben daher kommt es, dass häufig solche Lehrsätze, von denen anfangs ein Beweis Jahre lang vergeblich gesucht war, späterhin sich auf mehreren ganz verschiedenen Wegen beweisen lassen. Sobald ein neuer Lehrsatz durch Induction entdeckt ist, hat man die Auffindung \emph{irgend eines} Beweises freilich als das erste Erforderniss zu betrachten: allein nachdem ein solcher geglückt ist, darf man in der höhern Arithmetik die Untersuchung nicht immer als abgeschlossen und die Aufspürung anderer Beweise als überflüssigen Luxus ansehen. Denn theils kommt man gewöhnlich auf die schönsten und einfachsten
\SourceLeaf{170}%
% A544-P0170-D01
Beweise nicht zuerst, und dann ist gerade die Einsicht in die wunderbare Verkettung der Wahrheit der höhern Arithmetik dasjenige, was einen Hauptreiz dieses Studiums ausmacht, und nicht selten wiederum zur Entdeckung neuer Wahrheiten führt. Aus diesen Gründen ist hier die Auffindung neuer Beweise für schon bekannte Wahrheiten öfters für wenigstens eben so wichtig anzusehen, als die Entdeckung der Wahrheit selbst. Kennern der höhern Arithmetik sind diese Betrachtungen nicht neu; man weiss, dass ein grosser Theil von \textsc{Eulers} Verdiensten um dieselbe in der Auffindung von Beweisen für Lehrsätze besteht, die schon von \textsc{Fermat} wie es scheint durch Induction gefunden waren.

% A544-P0170-D02
Die Lehre von den quadratischen Resten gibt einen einleuchtenden Beleg zu dem vorhin Gesagten. Sie beruhet hauptsächlich auf dem sogenannten Fundamental-Theorem, welches darin besteht, dass die wechselseitigen Relationen zweier (ungeraden positiven) Primzahlen zu einander (in sofern der eine quadratischer Rest oder Nichtrest der andern ist) einerlei sind, so oft eine der Primzahlen oder beide unter der Form $4k+1$ stehen, entgegengesetzt aber, so oft beide Primzahlen von der Form $4k+3$ sind. Für solche Leser, die mit der höhern Arithmetik weniger bekannt sind, erinnern wir, dass eine ganze Zahl quadratischer Rest einer andern heisst, wenn die erstere um ein Vielfaches der andern vermehrt ein Quadrat geben kann; Nichtrest hingegen, wenn dies nicht möglich ist. Die Geschichte dieses schönen durch Induction äusserst leicht zu findenden Lehrsatzes wollen wir hier nicht vollständig wiederholen, sondern nur bemerken, dass der Verfasser vorliegender Abhandlung, nach Anfangs ziemlich lange vergeblich angestellten Untersuchungen, nach und nach bereits vier unter sich ganz verschiedene Beweise gegeben hat, wovon zwei in den \emph{Disquisitionibus Arithmeticis} enthalten sind, der dritte den Gegenstand einer eigenen Abhandlung im sechzehnten Bande der Commentationen ausmacht, und der vierte in eine Abhandlung \emph{summatio quarumdam serierum singularium} im ersten Bande der \emph{Commentationes recentiores} verwebt ist; über diese beiden Abhandlungen kann man unsere Anzeigen 1808. Mai 12 und Sept. 19 nachsehen, wo auch vollständigere geschichtliche Nachweisungen befindlich sind. Dass der Verf. bei diesen vier Beweisen, ungeachtet jeder derselben für sich in Rücksicht auf Strenge nichts zu wünschen übrig lässt, noch nicht stehen geblieben ist, bedarf zwar bei den Freunden der höhern Arithmetik keiner Rechtfertigung; indessen würde er doch wahrscheinlich sich nicht so eifrig bemüht haben, jenen Beweisen noch andere hinzuzufügen, wenn
